% =========================================================
% PREAMBLES.TEX - FINAL REPLACEMENT
% Untuk Laporan Kerja Praktik + Gambar Teknis PV
% Compile menggunakan XeLaTeX atau LuaLaTeX
% =========================================================

% --- PENGATURAN WARNA (HARUS DI ATAS) ---
% Memanggil xcolor di awal mencegah bentrok opsi dengan paket lain
\usepackage[table,xcdraw]{xcolor}

% --- PENGATURAN BAHASA DAN FONT ---
\usepackage[indonesian]{babel}

% Menggunakan font asli Times New Roman
% WAJIB compile dengan XeLaTeX atau LuaLaTeX, bukan pdfLaTeX
\usepackage{fontspec}
\setmainfont{Times New Roman}

% --- PENGATURAN SPASI DAN PARAGRAF ---
\usepackage{setspace}
\onehalfspacing

\usepackage{indentfirst}
\setlength{\parindent}{1cm}
\setlength{\parskip}{0pt}

\tolerance=1
\emergencystretch=\maxdimen
\hyphenpenalty=10000
\hbadness=10000

% --- PENGATURAN HALAMAN DAN MARGIN ---
\usepackage{geometry}
\geometry{
    top=3.5cm,
    left=3.75cm,
    bottom=2.5cm,
    right=2.75cm,
}

\usepackage{scrextend}
\usepackage{fancyhdr}
\usepackage{changepage}
\usepackage{pdflscape}
\usepackage{rotating}

% --- PENGATURAN LEVEL DAN PENOMORAN BAB/SUBBAB ---
\setcounter{secnumdepth}{6}

\renewcommand{\thesection}{\arabic{chapter}.\arabic{section}\hspace{0.05cm}}
\renewcommand{\thesubsection}{\arabic{chapter}.\arabic{section}.\arabic{subsection}\hspace{-0.25cm}}
\renewcommand{\thesubsubsection}{\Alph{subsubsection}.\hspace{-0.25cm}}
\renewcommand{\theparagraph}{\arabic{paragraph}.\hspace{-0.25cm}}

% --- PENGATURAN JUDUL BAB DAN SUBBAB ---
\usepackage{titlesec}

\titleformat{\chapter}
  {\doublespacing\fontsize{14pt}{16pt}\bfseries}
  {\MakeUppercase{\chaptertitlename\ \Roman{chapter}}\filcenter}
  {0.15cm}
  {\centering\uppercase}

\titlespacing*{\chapter}{0pt}{-1cm}{20pt}

\titleformat*{\section}{\fontsize{12pt}{14pt}\selectfont\bfseries}
\titleformat*{\subsection}{\fontsize{12pt}{14pt}\selectfont\bfseries}
\titleformat*{\subsubsection}{\fontsize{12pt}{14pt}\selectfont\bfseries}
\titleformat*{\paragraph}{\normalfont\fontsize{12pt}{14pt}\selectfont}

\titlespacing*{\section}{0pt}{10pt}{0cm}
\titlespacing*{\subsection}{0pt}{10pt}{0cm}
\titlespacing*{\subsubsection}{0pt}{10pt}{0cm}
\titlespacing*{\paragraph}{0pt}{10pt}{0cm}

% --- PENGATURAN DAFTAR ISI, GAMBAR, DAN TABEL ---
\usepackage[subfigure]{tocloft}

\renewcommand{\cfttoctitlefont}{\hfil\large\bfseries\MakeUppercase}
\renewcommand{\cftloftitlefont}{\hfil\large\bfseries\MakeUppercase}
\renewcommand{\cftlottitlefont}{\hfil\large\bfseries\MakeUppercase}

\renewcommand\cftchappresnum{BAB }
\renewcommand\cftchapaftersnum{}

\newlength\mylen
\settowidth\mylen{\bfseries BAB 1 :\ }
\cftsetindents{chap}{0pt}{\mylen}

\newcommand*{\noaddvspace}{\renewcommand*{\addvspace}[1]{}}
\addtocontents{lof}{\protect\noaddvspace}

% --- PENGATURAN GAMBAR DAN CAPTION ---
\usepackage{graphicx}
\graphicspath{{gambar/}}

\usepackage{float}
\usepackage{wrapfig}
\usepackage[hang,centerlast,nooneline,small,md]{subfigure}
\usepackage[font=footnotesize,format=plain,labelfont=bf,up,textfont=up]{caption}

\captionsetup[figure]{labelsep=space}
\captionsetup[table]{labelsep=space}

\renewcommand{\thefigure}{\arabic{chapter}.\arabic{figure}}
\renewcommand{\thetable}{\arabic{chapter}.\arabic{table}}

\usepackage{url}

\newcommand*{\putsource}[1]{%
    \textbf{\small{Sumber: }} \small{\url{#1}}
}

% --- PENGATURAN TABEL ---
\usepackage{array}
\usepackage{booktabs}
\usepackage{longtable}
\usepackage{tabularx}
\usepackage{multirow}
\usepackage{siunitx}

\sisetup{
    detect-all,
    group-separator={,},
    group-minimum-digits=4,
    round-mode=places,
    round-precision=3
}

% --- PENGATURAN MATEMATIKA ---
\usepackage{amsmath}
\usepackage{amssymb}
\usepackage{amsfonts}
\usepackage{stmaryrd}
\usepackage{mathtools}

\renewcommand{\theequation}{\arabic{chapter}.\arabic{equation}}

% --- PENGATURAN TIKZ, CIRCUITIKZ, DAN PGFPLOTS ---
% Bagian ini penting untuk gambar teknis pada Bab 3 dan Bab 4
\usepackage{tikz}
\usepackage{tikz-cd}
\usepackage{circuitikz}
\usepackage{pgfplots}

\usetikzlibrary{
    arrows.meta,
    positioning,
    calc,
    shapes.geometric,
    decorations.pathreplacing,
    matrix,
    patterns
}

\pgfplotsset{compat=1.18}

% --- STYLE TIKZ UNTUK GAMBAR TEKNIS PV ---
% Style ini dibutuhkan oleh gambar dari paper teknis:
% pvmod, pvmodbad, rblock, lab, labred, arr
\tikzset{
    pvmod/.style={
        draw,
        rectangle,
        rounded corners=2pt,
        minimum width=1.2cm,
        minimum height=0.6cm,
        line width=0.6pt,
        fill=blue!8,
        font=\scriptsize
    },
    pvmodbad/.style={
        draw,
        rectangle,
        rounded corners=2pt,
        minimum width=1.2cm,
        minimum height=0.6cm,
        line width=0.6pt,
        fill=red!15,
        font=\scriptsize
    },
    rblock/.style={
        draw,
        rectangle,
        minimum width=0.9cm,
        minimum height=0.35cm,
        line width=0.6pt,
        fill=white,
        font=\scriptsize
    },
    lab/.style={
        font=\scriptsize,
        align=center
    },
    labred/.style={
        font=\scriptsize,
        align=center,
        text=red!70!black
    },
    arr/.style={
        -{Latex[length=2mm]},
        line width=0.8pt
    }
}

% --- PENGATURAN LIST (ENUMERATE/ITEMIZE) ---
\usepackage{enumitem}

\newenvironment{packed_enum}{
    \begin{enumerate}[leftmargin=1.5\parindent]
        \setlength{\itemsep}{0pt}
        \setlength{\parskip}{0pt}
        \setlength{\parsep}{0pt}
}{
    \end{enumerate}
}

\newenvironment{packed_item}{
    \begin{itemize}[leftmargin=1.5\parindent]
        \setlength{\itemsep}{0pt}
        \setlength{\parskip}{0pt}
        \setlength{\parsep}{0pt}
}{
    \end{itemize}
}

% --- PENGATURAN REFERENSI DAN SITASI ---
\usepackage{cite}
\usepackage{varioref}
\usepackage[hidelinks]{hyperref}
\usepackage{cleveref}

% --- PENGATURAN LAIN-LAIN ---
\usepackage{blindtext}
\usepackage{lipsum}
\usepackage{epigraph}

% --- PENGATURAN LISTING KODE ---
\usepackage{listings}

% Warna Kustom untuk Listing
\definecolor{mygreen}{rgb}{0,0.6,0}
\definecolor{mygray}{rgb}{0.5,0.5,0.5}
\definecolor{mymauve}{rgb}{0.58,0,0.82}
\definecolor{codegreen}{rgb}{0,0.6,0}
\definecolor{codegray}{rgb}{0.5,0.5,0.5}
\definecolor{codepurple}{rgb}{0.58,0,0.82}
\definecolor{backcolour}{rgb}{0.95,0.95,0.92}

% Warna Khusus Processing
\definecolor{processingblack}{RGB}{0,0,0}
\definecolor{processinggray}{RGB}{102,102,102}
\definecolor{function}{RGB}{0,102,153}
\definecolor{lightgreen}{RGB}{102,153,0}
\definecolor{bluegreen}{RGB}{51,153,126}
\definecolor{magenta}{RGB}{217,74,122}
\definecolor{orange}{RGB}{226,102,26}
\definecolor{purple}{RGB}{125,71,147}
\definecolor{green}{RGB}{113,138,98}

% Style Default Listing
\lstdefinestyle{custom}{
    frame=single,
    backgroundcolor=\color{backcolour},
    commentstyle=\color{codegreen},
    keywordstyle=\color{magenta},
    numberstyle=\tiny\color{codegray},
    stringstyle=\color{codepurple},
    basicstyle=\ttfamily\footnotesize,
    breakatwhitespace=false,
    breaklines=true,
    captionpos=b,
    keepspaces=true,
    numbers=left,
    numbersep=6pt,
    showspaces=false,
    showstringspaces=false,
    showtabs=false,
    tabsize=2
}

\lstset{style=custom}

% Definisi sederhana bahasa Processing
% Definisi ini dibuat ringkas agar tidak memberatkan preamble.
% Jika laporan tidak memakai kode Processing, bagian ini tidak masalah dibiarkan.
\lstdefinelanguage{Processing}{
    morekeywords={
        abstract,break,class,continue,default,enum,extends,false,final,finally,
        implements,import,instanceof,interface,native,new,null,package,private,
        protected,public,static,strictfp,throws,transient,true,void,volatile,
        return,super,this,throw,catch,do,for,if,else,switch,synchronized,while,try,
        width,height,frameCount,frameRate,key,keyCode,keyPressed,mouseButton,
        mousePressed,mouseX,mouseY,pmouseX,pmouseY,pixels,
        Array,ArrayList,Boolean,Byte,Character,Class,Double,Float,Integer,
        HashMap,PrintWriter,String,StringBuffer,StringBuilder,Thread,
        boolean,byte,char,color,double,float,int,long,short,
        setup,draw,keyReleased,keyTyped,mouseClicked,mouseDragged,mouseMoved,
        mouseReleased,mouseWheel,settings,
        abs,acos,alpha,ambient,ambientLight,append,applyMatrix,arc,arrayCopy,
        asin,atan,atan2,background,beginCamera,beginContour,beginRaw,beginRecord,
        beginShape,bezier,bezierDetail,bezierPoint,bezierTangent,bezierVertex,
        binary,blend,blendColor,blendMode,blue,box,brightness,camera,ceil,
        circle,clear,clip,colorMode,concat,constrain,copy,cos,createFont,
        createGraphics,createImage,createInput,createOutput,createReader,
        createShape,createWriter,cursor,curve,curveDetail,curvePoint,
        curveTangent,curveTightness,curveVertex,day,degrees,delay,
        directionalLight,displayDensity,dist,ellipse,ellipseMode,emissive,
        endCamera,endContour,endRaw,endRecord,endShape,exit,exp,expand,fill,
        filter,floor,frustum,fullScreen,get,green,hex,hint,hour,hue,image,
        imageMode,join,launch,lerp,lerpColor,lightFalloff,lights,line,
        loadBytes,loadFont,loadImage,loadJSONArray,loadJSONObject,loadPixels,
        loadShader,loadShape,loadStrings,loadTable,loadXML,log,loop,mag,map,
        match,matchAll,max,millis,min,minute,month,nf,nfc,nfp,nfs,noClip,
        noCursor,noFill,noise,noiseDetail,noiseSeed,noLights,noLoop,norm,
        normal,noSmooth,noStroke,noTint,ortho,perspective,pixelDensity,point,
        pointLight,popMatrix,popStyle,pow,print,printArray,println,printMatrix,
        pushMatrix,pushStyle,quad,quadraticVertex,radians,random,randomGaussian,
        randomSeed,rect,rectMode,red,redraw,requestImage,resetMatrix,reverse,
        rotate,rotateX,rotateY,rotateZ,round,saturation,save,saveFrame,scale,
        screenX,screenY,screenZ,second,selectFolder,selectInput,selectOutput,
        set,shader,shape,shapeMode,shearX,shearY,shininess,shorten,sin,size,
        smooth,sort,specular,sphere,sphereDetail,splice,split,splitTokens,
        spotLight,sq,sqrt,square,stroke,strokeCap,strokeJoin,strokeWeight,
        subset,tan,text,textAlign,textAscent,textDescent,textFont,textLeading,
        textMode,textSize,texture,textureMode,textureWrap,textWidth,thread,
        tint,translate,triangle,trim,unbinary,unhex,updatePixels,vertex,year
    },
    sensitive=true,
    morecomment=[l][\color{processinggray}]{//},
    morecomment=[s][\color{processinggray}]{/*}{*/},
    morecomment=[s][\color{processinggray}]{/**}{*/},
    morestring=[b][\color{purple}]",
    morestring=[b][\color{purple}]',
    keywordstyle=\color{function}\bfseries,
    commentstyle=\color{processinggray},
    stringstyle=\color{purple}
}

% =========================================================
% AKHIR PREAMBLES.TEX
% =========================================================